\subsection{Standard Number Sets} 
There are certain sets of numbers that we talk about regularly in algebra.
\begin{center}
\begin{tabularx}{\textwidth-1pt}{|l|>{\hsize=0.375in\centering\arraybackslash}X|>{\hsize=0.375in\centering\arraybackslash}X|X|}
\hline
\textbf{Name}&\multicolumn{2}{|c|}{\textbf{Symbol}}&\textbf{Description}\\\hline
Natural Numbers&$\mathbb{N}$&&The counting numbers 1, 2, 3, ...\\\hline
Integers&$\mathbb{Z}$&&The natural numbers, their opposites $-1$, $-2$, $-3$, ..., and 0. Numbers with no fractional part.\\\hline
Rational Numbers&$\mathbb{Q}$&&The numbers that can be represented as fractions $\frac{p}{q}$ where $p$ and $q$ are integers, $q\neq 0$.\\\hline
Real Numbers&$\mathbb{R}$&&The rational numbers, together with irrational numbers to ``fill in'' the number line.\\\hline
\end{tabularx}
\end{center}
\begin{note}We will sometimes refer to the \term{positive integers} to mean the natural numbers, and the \term{nonnegative integers} to mean the natural numbers plus zero.\end{note}
\begin{note}The same rational number can be written in many forms, e.g. \mbox{$\frac{3}{5}=\frac{6}{10}=\frac{-9}{-15}$} etc. These are all the same rational number.\end{note}
\begin{example}
For each number below, determine which of the standard number sets contains it.
\begin{lpic}{}{7}
\foreach \x/\xlbl in {3/N,5/Z,7/Q,9/R}
	\regtext{\x+.5}{-1}{$\mathbb{\xlbl}$};
\regtextleft{11}{-1}{Notes};
\foreach \y/\ylbl in {1/6,2/-\frac{16}{4},3/\frac{3}{5},4/-0.78,5/0.\overline{3},6/\sqrt{2}}
	\regtextright{2}{-1-\y}{$\ylbl$};
\end{lpic}
\end{example}
How can we recognize rational and irrational numbers?
\begin{itemize}
\item All \makeblank are rational.
\item Any numbers written as a \makeblank is rational.
\item Every \makeblank[1.5in] or \makeblank[1.5in] decimal is rational.
\item All other decimals are irrational.
\item Square roots of positive integers are either \makeblank or irrational.
\end{itemize}